\chapter{Клеточные автоматы, обратимость и представление Маделунга}
\label{app:background}

\section{Зачем это приложение}

В основном тексте мы говорим о микроуровне~$M$, локальном законе~$g$ и переходе к макроописанию~$T$,
не ожидая от читателя специализации по теории клеточных автоматов или по истории квантовой гидродинамики.
Тем не менее несколько терминов --- «клеточный автомат'', «обратимый шаг'', «представление Маделунга'' ---
встречаются в гл.~\ref{ch:axiom-rationale}--\ref{ch:matter} и заслуживают краткого,
самодостаточного пояснения.
Здесь --- не полный учебник, а мост между классической литературой и нашей моделью.

\section{Что такое клеточный автомат}

\emph{Клеточный автомат} --- дискретная динамическая система на решётке.
Пространство состояния --- множество конфигураций $\{s(x)\}_{x\in\Lambda}$,
где $x$ --- узел решётки, $s(x)$ --- состояние в узле (конечное или счётное множество значений).
Время дискретно: $t=0,1,2,\ldots$
За один \emph{такт} каждый узел обновляется \emph{локально}:
новое значение $s(x,t+1)$ зависит только от значений в некоторой окрестности $N(x)$ и,
возможно, от собственного прошлого значения.
Общий вид --- $s(x,t+1)=f\bigl(\{s(y,t)\}_{y\in N(x)}\bigr)$.

Идея восходит к работам \emph{Станислава Улама} и \emph{Джона фон Неймана} (1940-е):
моделировать сложное поведение из простых локальных правил на решётке.
Классический пример --- \emph{«Игра жизни''} Конвея (1970): двумерная решётка,
два состояния «живая/мёртвая'', правило рождения и гибели по числу соседей.
Она иллюстрирует \emph{универсальность} клеточных автоматов --- способность воспроизводить
произвольные вычисления --- но не претендует на описание физического вакуума.

В нашей монографии микроуровень~$M$ \emph{математически} есть клеточный автомат:
счётная решётка~$\Lambda$ (FCC), дискретное время~$\htick$, локальное правило~$g$.
Отличие от «игровых'' автоматов --- физические ограничения $A_1$--$A_{16}$:
унитарность, термодинамика, топология заряда, спин~$1/2$, а не произвольная таблица переходов.
Закон~$g$ не подбирается для красивой картинки, а \emph{замыкается} из этих условий
(гл.~\ref{ch:evolution}).

\subsection{Окрестности: фон Нейман, Мур, световой конус}

Часто различают \emph{окрестность фон Неймана} (общие грани, без диагоналей)
и \emph{окрестность Мура} (все узлы в квадрате, включая диагонали).
Условие~$A_1$ в $(3{+}1)$ измерениях задает не произвольный Мур или $N_4$,
а \emph{светоподобную} первую оболочку: за один такт~$\htick$ сигнал может пройти
только на соседа на расстоянии~$\hl$ по световому конусу (гл.~\ref{ch:carrier}, \cref{def:epsilon}).
«Диагональ за такт'' --- запрещённый сверхсветовой шаблон, не техническая мелочь.

\section{Тьюринг-полнота и зачем вообще КА-подход}
\label{sec:app-universality}

\subsection{Вычислительная универсальность}

\emph{Тьюринг-полнота} (вычислительная универсальность) означает:
система может имитировать машину Тьюринга с произвольной программой и входом,
если ей выделено достаточно места и времени.
Не «умность'' правила, а \emph{полнота класса} динамик: локальный автомат на решётке
способен кодировать произвольные конечные логические схемы и последовательности операций.

Исторически:
\begin{itemize}
\item \textbf{фон Нейман} (1950-е) --- самовоспроизводящийся автомат:
  доказательство принципа, что локальное правило на решётке может нести
  произвольную структурированную информацию (книга 1966, посмертно);
\item \textbf{Конвей} --- «Игра жизни'': универсальность на $(2{+}1)$-решётке с двумя состояниями;
\item \textbf{Wolfram, Rule~110} (строгое доказательство Cook, 2004) ---
  даже одномерный элементарный автомат может быть Тьюринг-полным;
\item \textbf{Тоффоли} --- обратимые CA и «бильярдная'' модель вычислений:
  универсальность совместима с \emph{обратимостью} каждого шага.
\end{itemize}

Отсюда следует не философский тезис «Вселенная --- ноутбук'',
а \emph{методологический}: если микрофизика --- клеточный автомат с конечным алфавитом
на счётной решётке, то \emph{класс} таких систем уже достаточен,
чтобы порождать сколь угодно сложную детерминированную динамику ---
включая ту, которая на~$T$ выглядит как квантовая механика и поля Стандартной модели.
КА-подход \emph{применим} не потому, что мы «выбрали игру'',
а потому что Тьюринг-полнота гарантирует: локальные правила не обречены быть тривиальными.

\subsection{Что из этого следует для нашей модели}

Микроуровень~$M$ с законом~$g$ --- \emph{конкретный} автомат, не произвольная таблица Rule~110.
Он не «программируется'' под каждый эксперимент: $g$ \emph{замыкается} из $A_1$--$A_{16}$
(гл.~\ref{ch:evolution}).
Универсальность CA объясняет, \emph{почему} такой каркас вообще годится для физики;
абсолютные условия объясняют, \emph{какой именно} автомат выделяется.

\begin{remark}[Три уровня, которые не путать]
\begin{enumerate}
\item \textbf{Класс CA} --- обладает Тьюринг-полнотой; локальность + дискретность $\Rightarrow$ богатая динамика.
\item \textbf{Закон~$g$} --- один элемент класса, физически ограниченный; обратим ($A_{13}$),
  унитарен ($A_3$), с FCC-носителем ($A_1$, гл.~\ref{ch:carrier}).
\item \textbf{Макроуровень~$T$} --- осреднение и статистика; «случайность'' и Born ---
  не универсальная программа в~$g$, а грубый прибор (гл.~\ref{ch:macro}, \cref{prop:determinism}).
\end{enumerate}
\end{remark}

Обратимая универсальность (Тоффоли, Фредкин) особенно близка нашей постановке:
физический шаг не должен стирать информацию, иначе унитарность~$A_3$ нарушена.
Тьюринг-полнота без обратимости --- возможна (Игра жизни необратима на многих конфигурациях);
наш~$g$ --- в пересечении: \emph{достаточно богат} и \emph{строго обратим}.

\section{Обратимость: Тоффоли, Фредкин и условие \texorpdfstring{$A_{13}$}{A13}}

\subsection{Зачем физике обратимый шаг}

На микроуровне~$M$ мы требуем, чтобы эволюция сохраняла информацию (условие~$A_3$)
и была \emph{обратимой}: существует $g^{-1}$ такое, что $g^{-1}(g(\Psi))=\Psi$
(условие~$A_{13}$).
Иначе часть микросостояний «схлопывалась'' бы в один образ --- необратимая диссипация
на фундаментальном шаге, несовместимая с унитарностью.

\subsection{Тоффоли и Фредкин}

В 1970--1980-е \emph{Эдвард Фредкин} и \emph{Томмазо Тоффоли}
(часто вместе с \emph{Норманом Марголусом}) развивали идею, что
\emph{вычисление и физика} могут быть дискретными и \emph{обратимыми}:
каждый шаг эволюции --- перестановка конечного множества состояний.

\begin{itemize}
\item \textbf{Фредкин} --- «цифровая физика'': Вселенная как большой автомат;
  известен \emph{вентиль Фредкина} --- универсальный обратимый логический элемент.
\item \textbf{Тоффоли} --- строгая теория \emph{обратимых} клеточных автоматов:
  если правило локально инъективно на конфигурациях с конечным алфавитом,
  оно глобально обратимо (теорема типа Garden of Eden / Curtis-Hedlund для обратимых систем).
  Книга Toffoli--Margolus \emph{Cellular Automata Machines} (1987) --- классический источник.
\end{itemize}

Их мотивация была в основном \emph{вычислительной и философской}
(обратимые вычисления, отсутствие энтропийного «стирания'' бита при необратимой логике AND).
В нашей модели тот же \emph{математический} принцип стоит на \emph{физическом} месте:
$A_3$ + $A_{13}$ $\Rightarrow$ микрошаг --- унитарная перестановка на $Z_N[i]$;
схема «лягушка'' с учётом фазового импульса и явным $g^{-1}$ (гл.~\ref{ch:evolution}).

\subsection{Чем мы не занимаемся}

Мы \emph{не} утверждаем, что Вселенная --- «компьютер Фредкина'' в буквальном смысле,
и не сводим всю физику к булевой логике вентилей.
Обратимость --- необходимое свойство \emph{локального закона}~$g$ при сохранении $N=\sum|z|^2$,
а не метафора программирования.
Статистика и необратимость \emph{измерения} возникают на~$T$ (гл.~\ref{ch:macro}),
когда осреднение~$\mathcal{B}$ перестаёт быть инъективным.

\section{Представление Маделунга}

\subsection{Исторический контекст}

В 1927 году \emph{Эрвин Маделунг} показал, что уравнение Шрёдингера для одной частицы
можно переписать в гидродинамической форме.
Пусть $\psi=\sqrt{\rho}\,e^{iS/\hbar}$ --- волновая функция.
Выделим \emph{плотность} $\rho=|\psi|^2$ и \emph{фазу} $S$.
Подстановка в Шрёдингера даёт два уравнения, формально похожие на гидродинамику:

\begin{enumerate}
\item \textbf{Уравнение неразрывности:}
$\partial_t\rho + \nabla\cdot(\rho\mathbf{v})=0$,
где $\mathbf{v}=\nabla S/m$ --- скорость «жидкости''.
\item \textbf{Уравнение импульса} (типа Эйлера) с дополнительным
\emph{квантовым давлением}, пропорциональным $\nabla^2\sqrt{\rho}/\sqrt{\rho}$.
\end{enumerate}

Так возник образ \emph{квантовой жидкости}: $\rho$ и $\mathbf{j}=\rho\mathbf{v}$
--- поля плотности и тока, как в классической гидродинамике,
плюс «квантовый'' член, исчезающий в классическом пределе.

\subsection{Маделунг в нашей монографии}

Маделung изначально --- язык \emph{непрерывного} $\mathbb{R}^3$ и дифференциальных уравнений.
В нашей модели он живёт на \emph{макроуровне}~$T$, после осреднения~$\mathcal{B}$
(гл.~\ref{ch:macro}, \ref{ch:matter}):

\begin{itemize}
\item на~$M$ нет гладкого $\partial_t\rho+\nabla\cdot\mathbf{j}=0$ как фундаментального закона;
  там --- унитарность~$A_3$ и дискретные учёты импульса на связях решётки;
\item на~$T$ осреднённое поле $\Phi=\mathcal{B}z$ допускает \emph{диагностику}
  в терминах Маделунга: $\rho\approx|\Phi|^2$, ток через фазовые градиенты,
  неразрывность как следствие осреднения, а не аксиома на~$M$.
\end{itemize}

Имя «решёточная жидкость'' в заглавии монографии отражает именно эту линию:
микроуровень --- несжимаемая (в смысле~$A_3$ и модуля~$K_P$) динамика на решётке;
макроуровень --- гидродинamika и уравнение Шрёдингера как \emph{эффективные} теории,
связанные с Маделунгом так же, как Navier--Stokes связан с кинетической теорией.

\begin{remark}[Частые заблуждения]
\begin{enumerate}
\item \textbf{«Маделунг доказывает дискретность.''} Нет: он переписывает \emph{уже} непрерывное уравнение Шрёдингера.
  Дискретность~$M$ выводится из $A_1$--$A_{16}$ (гл.~\ref{ch:foundations}).
\item \textbf{«На M должна выполняться неразрывность Маделунга.''} Нет: остаток на одном такте
  ожидаем; гладкая неразрывность --- свойство~$T$ при $\lambda\gg\hl$.
\item \textbf{«Клеточный автомат = Игра жизни.''} Игра жизни --- частный пример;
  наш~$g$ --- обратимый спинорный автомат с фазой на $Z_{512}[i]$, не булевому автомату.
\end{enumerate}
\end{remark}

\section{Сводная таблица}

\begin{table}[ht]
\centering
\caption{От классических понятий к главам монографии}
\label{tab:app-map}
\small
\begin{tabular}{@{}p{3.2cm}p{4.2cm}p{5.8cm}@{}}
\toprule
Понятие & Классика & У нас \\
\midrule
Клеточный автомат & Улама, von Neumann; Wolfram & $M$, закон~$g$, гл.~\ref{ch:evolution} \\
Тьюринг-полнота & Конвей, Rule~110; Toffoli & Класс CA; см. \S\ref{sec:app-universality} \\
Обратимость & Фредкин, Тоффоли, Марголус & $A_{13}$, схема «лягушка'', $g^{-1}$ \\
Окрестность & Moore, von Neumann & Световой конус, $A_1$, \cref{def:epsilon} \\
Маделунг & 1927, $\rho$, $\mathbf{j}$, квантовое давление & Только~$T$, гл.~\ref{ch:macro}, \ref{ch:matter} \\
\bottomrule
\end{tabular}
\end{table}

\section{Литература для первого чтения}

\begin{itemize}
\item E.~Fredkin, Digital mechanics (обратимая дискретная физика, 1990-е).
\item T.~Toffoli, N.~Margolus, \emph{Cellular Automata Machines} (MIT Press, 1987).
\item J.~von Neumann, \emph{Theory of Self-Reproducing Automata} (1966).
\item E.~Маделунг, Quantentheorie in hydrodynamischer Form (1927).
\item M.~Cook, Universality in elementary cellular automata (2004) --- Rule~110.
\item Обзор CA: S.~Wolfram, \emph{A New Kind of Science} (2002) --- с оговоркой:
  наша модель не следует его классификации «элементарных'' правил; физика задаёт~$g$.
\end{itemize}

\medskip
\noindent
В основном тексте достаточно помнить: \textbf{M} --- обратимый локальный автомат на решётке,
достаточный по классу (Тьюринг-полнота CA) и конкретный по закону~$g$;
\textbf{T} --- макроописание, где уместны Маделунг, волновые уравнения и SI.
